\chapter{INTRODUCTION}
\section{Background }
Malaria remains a major public health burden, with an estimated 282 million cases reported in 2024, out of which 95\% occurred in the Sub-Saharan part of Africa \cite{WHO_malaria_2024}. Despite the implementation of various control strategies, including usage of insecticide-treated nets(ITNs), indoor residual spraying(IRS) and effective treatment, malaria continues to cause a high number of illnesses and deaths each year \cite{fullman2013nets}. In 2024 alone, the World Health Organization(WHO) African reported approximately 579,000 deaths ; out of which 80\% are children under five years \cite{WHO_malaria_2024}. This shows how hard it is to manage malaria in our current world. One of the key challenges in malaria control is that the disease is not evenly distributed across space and time. Some areas experience consistently high transmission, while others face seasonal outbreaks \cite{amek2012spatial}. This spatial and temporal variability makes malaria difficult to control using uniform intervention strategies. Because of this nature of the disease, public health systems require  accurate prediction that can show where and when malaria risks are very high so that they can apply intervention strategies \cite{Mategula2025}.

Malaria transmission is influenced by multiple interacting factors that vary across locations and over time \cite{Odhiambo2020,amek2012spatial}. Climatic factors like rainfall, temperature, and humidity, plays an important role in influencing mosquito breeding, parasite growth, and vector survival \cite{Caminade2023}.  Environmental factors like vegetation, land cover, and elevation influence mosquito breeding habitats \cite{Kuenzer2022}. Human migration, population density and access to health facilities have an impact on malaria transmission \cite{Chirebvu2014}. The relationships between these various elements are complicated and nonlinear \cite{Okundalaye2025Climate-based}. For example, moderate rainfall promotes mosquito breeding while excessive rainfall washes away breeding places for mosquitoes. Again, changes in climatic conditions do not lead to immediate changes in malaria cases due to lag effects \cite{Laneri2010}. Such  kind of complexities make malaria prediction challenging, especially at  geographical scales such as counties and sub-counties, where public health decisions are made \cite{Okunlola2019}. Lack of effective malaria prediction systems may lead to misuse of the already limited public health resources, where efforts are focused on low-risk areas while ignoring the highly-risky areas altogether \cite{Mategula2025}. Bringing out the need for robust predictive models that can better capture the complex nature of malaria transmission and provide actionable insights for decision-makers.

Generally, the prediction and mapping of malaria have depended on statistical approaches such as regression models, time series analysis, and spatial modeling \cite{Held2022}. While these methods have contributed to understanding general malaria patterns, they have several limitations since they make assumptions that malaria patterns depend on simple linear relationships and are unable to incorporate non-linear interactions, temporal delays, or interactions(spatial) between adjacent regions in space \cite{Odhiambo2020}. In recent years, Machine learning(ML) approaches such as Random Forest(RF), XGBoost, and LightGBM have been employed in predicting malaria due to their capacity to work well with large data sets and capture complex nonlinear relationships \cite{Syafei2023, Yang2024}. These models are relatively good at making predictions, but they often fail to capture the spatio-temporal nature of malaria patterns, which is an essential characteristic of malaria transmission.

Recently, deep learning approaches have provided new opportunities for improving malaria prediction. Models such as Convolutional(CNNs), Long Short Term Memory(LSTMs), Gated Recurrent Units(GRUs), and Transformers have been applied in spatio-temporal prediction tasks due to their ability to capture complex spatial and temporal dependencies \cite{Rasheed2023}. When combined, both the spatial and temporal dependencies could be modelled effectively, making them suitable for representing complex malaria transmission dynamics. Nonetheless, applications of deep learning in the field of malaria  forecasting through  hybrid approaches remains relatively unexplored.
Previous studies have mainly focused on traditional statistical and machine learning techniques, but limited attention has been given to jointly modelling the spatio-temporal nature of malaria transmission. In addition, hybrid architectures such as CNN-LSTM, CNN-GRU, and CNN-Transformer have not been sufficiently compared under similar epidemiological and environmental conditions. As a result, there is still limited understanding of how these hybrid models perform in capturing heterogeneous spatio-temporal malaria transmission patterns across Kenyan counties. 
We address these limitations by proposing a spatio-temporal malaria prediction framework that integrates malaria incidence data with climatic, environmental, demographic, and intervention-related variables. The framework would combine hybrid deep learning approaches to better capture the complex spatial and temporal patterns of malaria transmission. Although malaria remains is still public health challenge across East African countries such as Kenya, Uganda, and Tanzania \cite{Snow2015}, We chose Kenya as the case study because because of its diverse ecological and climatic conditions \cite{amek2012spatial,Gething2023} in malaria transmission. To capture these differences, the study uses data from all 47 counties, enabling fine-scale spatio-temporal analysis and forecasting of malaria risk. The study therefore, purpose to predict  county-level malaria risk maps and early warning indicators to support targeted interventions and evidence-based public health decision-making.
\section{Statement of the Problem}
Malaria transmission is influenced by a combination of climatic, environmental, and human-related factors that vary across locations and over time \cite{amek2012spatial,Okunlola2019}. In Kenya, malaria risk is not evenly distributed across the country. Some counties experience consistently high transmission, while others face seasonal outbreaks due to differences in climatic, environmental, and population-related conditions \cite{Gething2023}. This spatial and temporal variation makes it difficult to accurately predict where and when malaria risk is likely to increase, particularly at the county level where decisions on disease control and resource allocation are often made \cite{Odhiambo2020}.

 Several approaches have been used to predict malaria risk, but many do not adequately capture the complex nature of malaria transmission. The relationships between malaria incidence and factors such as rainfall, temperature, vegetation, and population dynamics are often nonlinear and may occur with time delays \cite{Syafei2023,Laneri2010}. In addition, malaria patterns in one area can be influenced by neighbouring areas, yet this spatial dependence is not fully represented in many existing models \cite{Odhiambo2020}. These limitations can reduce the accuracy and practical usefulness of prediction results.

 Recent developments in machine learning and deep learning have improved predictive modelling; however, their application to malaria forecasting at county level in Kenya remains limited \cite{Odhiambo2020,Gething2023}. In particular, hybrid approaches that combine CNNs with temporal architectures such as LSTM, GRU, and Transformers have not been widely explored for modelling malaria transmission despite their ability to learn both spatial and temporal patterns \cite{Rasheed2023}.

 This study therefore develops and evaluates a hybrid CNN-based spatio-temporal framework for malaria prediction in Kenya. By integrating malaria incidence data with climatic, environmental, demographic, and intervention-related variables, the study compares CNN-LSTM, CNN-GRU, and CNN-Transformer models for county-level malaria risk prediction and mapping, with the aim of supporting early warning and malaria control planning.
\section{Research Objectives}

\subsection{Main Objective}
To develop and evaluate spatio-temporal machine learning and hybrid deep learning models for malaria risk prediction across Kenyan counties.
\subsection{Specific Objectives}

\begin{enumerate}
	\item To model spatial and temporal dependencies in malaria transmission using hybrid deep learning architectures.
	\item To integrate climatic, environmental, and demographic variables into a unified spatio-temporal prediction framework.
	\item To compare the predictive performance of baseline machine learning and hybrid deep learning architectures using RMSE, MAE, and $R^2$
	metrics.
	\item To interpret malaria prediction patterns using SHAP explainability analysis and spatial malaria risk mapping.
\end{enumerate}

\section{Research Questions}

\begin{enumerate}
	\item Can hybrid deep learning models accurately capture spatial and temporal dependencies in malaria transmission?
	\item How effectively can climatic, environmental, and demographic variables be integrated to improve malaria risk prediction?
	\item Do hybrid deep learning models  perform better than traditional machine learning approaches in predictive accuracy and hotspot detection?
	\item How can SHAP explainability analysis and high-resolution malaria risk  spatial mapping support interpretation of malaria prediction useful for public health decision-making?
\end{enumerate}
\section{Justification }
Many malaria predictive models  often perform poorly in real operational settings since they do not fully take into account spatial dependency,temporal changes, and nonlinear nature of malaria transmission. Many of these models also rely on simplified assumptions and provide predictions at coarse spatial scales, which limits their usefulness for targeted malaria interventions at local levels. Therefore a need for scalable modelling approaches that can be applied across all counties while still maintaining strong predictive performance. A spatio-temporal machine learning framework allows the integration of climatic, environmental, demographic, and intervention-related datasets to support county-level malaria prediction and spatial risk mapping. This can improve the identification of localized malaria risk patterns and support more targeted public health responses. In addition, integrating predictive models into health information systems like DHIS2 boosts the effectiveness of malaria control initiatives, promotes early warning systems, and facilitates real-time decision-making.

\section{Significance of the Study}
This study contributes to spatio-temporal malaria prediction by integrating climatic, environmental, and demographic factors using hybrid deep learning models. This could help in making decisions regarding the appropriate measures to undertake as well as allocation of malaria management resources in Kenya. In addition, the findings can offer useful information that can help develop effective malaria control strategies in counties and sub-counties.






